%!TEX root = main.tex
\documentclass[main.tex]{subfiles}

\begin{document}

\chapter{Conclusion}

\vspace{2em}

\setlength{\epigraphwidth}{0.65\textwidth}
\renewcommand{\textflush}{flushepinormal}
\epigraph{\corr{Have we addressed the transcription bottleneck?~The answer is a qualified yes.~[...] We are more selective in where to ``transcribe exhaustively.'' This is how we manage the transcription trap.~[...] [This approach] offers diverse workflows and flexibility in the face of diverse language situations.~Finally, it opens up new opportunities to engage with speakers.~These people offer the deepest knowledge about the language and the best chance of scaling up the work.\linebreak{}This is their language.}}%
         {\justifying{}\corr{Bird, S. (2020). Sparse transcription. \emph{Computational Linguistics}, 46(4), 713–744.}}

\vspace{2em}

\corr{The goal of language documentation is to produce a multi-purpose record of a language \parencite{himmelmann1998documentary}.~Given this goal, the recorded materials may encompass a variety of linguistic topics and cultural practices and, consequently, materials appropriate for one purpose (e.g.~linguistic analyses) may be entirely inappropriate for others (e.g.~creating language learning materials).~Thus, the ability to efficiently locate and extract only the relevant information from recorded materials is a need shared among in wide variety of use cases --- from linguistic and lexicographical analyses, to development of language materials, to simply wanting to listen to the speech of one's elders.}

\corr{Plain text searches on transcriptions of the spoken content have been the primary pathway for search and retrieval of material in documentation corpora.~Yet as recording speech is much easier than transcribing the recorded speech, many recordings collected for documentation efforts remain untranscribed.~This leads to a transcription bottleneck, which in turn limits how easily untranscribed materials can be accessed.~There are several contributing factors to the transcription bottleneck --- from the need for a writing system and transcribers literate in this system to inefficiencies in the conventional transcription workflow.}

\corr{Towards alleviating or bypassing this bottleneck altogether, I substantiated the importance of addressing two research questions:~first, how to bypass the need for transcription and derive a searchable index, and second, how to make most effective use of available resources to develop automatic speech recognition (ASR) to help accelerate the transcription work.~I argued that solutions that address each of these questions can be broadly viewed as diversifying the pathways that can help meet diverse information needs.~Leveraging recently developed pre-trained models for speech processing, I examined for each pathway two possible implementations that cater to a common but distinct scenario within language documentation projects.}

\corr{For scenarios where the speech corpora comprise recordings only in the target language(s), I examined the use of pre-trained models as part of a spoken term detection system.~Such systems permit users to perform voice-based searches using spoken queries and in turn permit confirmed hits to act as an index of the corpus \parencite{bird-2020-sparse,le-ferrand-etal-2020-enabling}.~Evaluated using real world documentation data from 8 languages, results showed that in the worst case scenario using pre-trained models on average can help retrieve 40-50\% of all occurrences of a given query.~Given that the alternative is scrubbing through untranscribed audio recordings in near real time, this method helps offer \emph{some} form of search and retrieval for uses requiring an exact match.~For more discovery-oriented uses (e.g.~finding near minimal pairs), however, this method can readily return the most similar sounding near matches (e.g.~\emph{halo}, \emph{hollow}, etc. given the query \emph{hello}).}

\corr{For scenarios where the speech corpora comprise recordings with a more widely spoken vehicular language used for commentary and questions (e.g.~\emph{What is the word for `soft'?}), I examined the use of pre-trained models to automatically isolate and transcribe the vehicular language.~The vehicular language transcriptions act as a searchable proxy index, indicating the approximate locations where certain target language words and phrases are being discussed.~Evaluated on a 136-hour corpus of Muruwari (an Australian Aboriginal language), results showed that pre-trained models can be readily adapted for this processing pipeline and reduce transcription time by 20\%.~The pipeline has helped meet information needs not only for anticipated uses such as finding words for developing an orthography but also novel, unanticipated uses such as mentions of a specific cultural artefact for use as evidence in repatriation proceedings.}

\corr{For scenarios where some external text sources are available, I examined whether leveraging such data can help kick-start the ASR development process when only limited amounts of transcribed speech is available (e.g.~10 minutes).~The text data can be used to develop a lexicon which limits ASR system output to possible words (e.g.~\emph{*dogz} $\rightarrow$ \emph{dogs}) or a language model which helps bias the ASR system to more probable word sequences (e.g. \emph{too dogs} $\rightarrow$ \emph{two dogs}), helping yield better quality transcriptions requiring fewer corrections.~This approach was evaluated using two pairs of minoritised languages: Gronings and Frisian (two West-Germanic languages spoken in the Netherlands), and Nasal and Besemah (two Malayo-Polynesian languages spoken in Indonesia).~Results showed that while this approached was promising for the two West-Germanic languages, typically over-represented in the data used to develop pre-trained models, ASR performance on the two under-represented Austronesian languages remained poor --- suggesting the need for a different approach for such languages.}

\corr{For scenarios where some untranscribed audio is available, I examined whether this data can be used to adapt pre-trained models to reduce the ASR performance gap resulting from under-representation.~In particular, I examined whether the effectiveness of this adaptation process could be increased through sourcing supplemental data from another similar language, i.e.~a donor language.~I hypothesised that acoustic-phonetic and phonotactic similarity between the target and donor languages predict the effectiveness (as measured by better ASR performance in the target language) and develop a novel measure to determine this similarity --- the Acoustic Token Distribution Similarity (ATDS).~Across a set of typologically different target languages (Punjabi, Galician, Iban, Setswana), I show that the ATDS between the target language and its candidate donors precisely predicts target language ASR performance and outperforms all other similarity measures examined.~In other words, this approach provides a means of using only untranscribed audio to optimise pre-trained models for developing ASR systems for a given target language, which in turn can help accelerate transcription.}

\section{Limitations and future directions}

\corr{To ensure findings were as informative as possible for applications in language documentation projects, I sought out data from real-world documentation projects for evaluating all implementations above --- except the last study on adapting pre-trained models.~In order to systematically study the effects of donor language choice on ASR performance in a target language, I used well-matched ASR datasets with careful read speech recorded in low-noise environments.~In practice, there may be many potential mismatches between the target language corpora (e.g.~read elicitations from a single speaker) and available donor language corpora (e.g.~noisy, multi-speaker conversations).~Further, given that the proposed similarity measure is an acoustic measure, its validity in comparing corpora with differing noise levels needs further investigation.~A follow-up experiment for example could examine the same target Punjabi data while sourcing acoustically mismatched donor data for Indic languages from the Shrutilipi dataset,\footnote{\url{https://ai4bharat.iitm.ac.in/shrutilipi/}} containing audio from radio archives.}

\corr{While the vehicular language isolation and transcription pipeline did examine data sourced from a language documentation project, the single-speaker self-elicitation diverges from the more common multi-speaker settings with typically one person eliciting material from one or more speakers of the target language.~Thus, the implementation of an equivalent pipeline will require the use of a speaker diarisation model, which helps differentiate who is speaking when.~Additionally, our timed experiments showed that using a cascaded approach where a voice activity detection model (to differentiate speech from silence) feeding into a separate spoken language identification model (to differentiate vehicular from target language speech) yielded poor segmentations that required extensive corrections.~Thus, for future implementations for multi-speaker recordings with a vehicular language, a promising approach for jointly performing segmentation and multi-class classification tasks is the recently proposed power-set approach \parencite{plaquet23_interspeech}, which for 2 speakers (S1, S2) and 2 languages (L1, L2) could help identify regions of S1L1, S1L2, S2L1, S2L2.}

\corr{Given their recency, the use of large pre-trained models comes with many practical infrastructure challenges which must be overcome before users can leverage them in their existing workflows on their own devices (e.g.~a laptop).~For example, for use as a feature extractor in a spoken term detection system, the extracted features of excerpts in the corpora would ideally be cached.~The efficient storage of these features, however, can require several hundred gigabytes of storage, depending on the size of the speech corpus.~Given such costs, all model training and subsequent processing of documentation data using the trained models were conducted by me on a university compute cluster, i.e. not on the machines of the collaborating linguists who collected the data.~Fortunately, given global commercial interests in reducing such costs, this current barrier is likely to decrease with time as usage of pre-trained models becomes increasingly ubiquitous.~Nevertheless, irrespective of such costs, how best to integrate available technologies into the daily workflow of end users remain a difficult, open question.~Some promising avenues towards finding answers are Elpis \parencite{foley2018building}, which offers a graphical user interface for ASR training (now including wav2vec 2.0 fine-tuning), and development of plug-ins for using pre-trained models from directly within Elan, e.g.~\texttt{SileroVAD-Elan} and \texttt{xls-r-elan}}.\footnote{Respectively: \url{https://github.com/l12maro/SileroVAD-Elan} and \url{https://github.com/coxchristopher/xls-r-elan}.}

\corr{Last, but not least, I acknowledge that all but one of the studies focused on evaluating system performance, which may or may not directly translate to meeting user needs.~On one hand, an implemented system may entirely suffice for meeting identified needs and even novel, unanticipated ones --- as in the case of finding mentions of a cultural artefact for use in repatriation proceedings.~On the other hand, there may be unanticipated barriers that prevent the adoption of a system, requiring design revisions.~For example, in piloting the deployment of a spoken term detection or `word spotting' system for Kunwinjku (an Australian Aboriginal language), \textcite{le2022learning} found that the very notion of what even  constitutes a word placed a barrier to understanding the purpose of the system by speakers of this polysynthetic language.~\textcite{le-ferrand-etal-2022-fashioning} redesigned the system to enable incremental and interactive `syllable' spotting, which is more naturally in line with how speakers decompose complex Kunwinjku utterances, and which speakers view as a tool for learning and practising literacy via on-screen orthographic representations of the spotted syllables --- a use that is in line with the community's priorities for language preservation.}

\section{A final note}

\corr{Linguists often find themselves explaining to a lay audience the difference between grammatical prescriptivism and descriptivism.~The first conjures up under-theorised notions of rigid practices while the second an openness and curiosity about the diverse ways in which language is used, how these ways vary between speakers and languages, and how these ways are constantly changing.~Reflecting on the seminal work of \textcite{himmelmann1998documentary}, for which we often cite the quintessential purpose of language documentation --- to create a multi-purpose record of a language ---, \textcite{himmelmann2018meeting} points out that the practice of transcription is severely under-theorised despite its central role in language documentation.~One effect of this under-theorisation appears to be certain rigid notions of the transcription process, assuming that all purposes require complete and exhaustive transcription.~This assumption and notions that underlie it have recently questioned by \textcite{bird-2020-sparse}, who argues that a key part of combating the transcription bottleneck is knowing when transcription is absolutely necessary and, crucially, when it isn't.~Of course, it will be no easy feat to establish this collective knowledge.~A first step, though, may simply be an openness and curiosity about the diverse ways in which spoken language recordings are used, how these uses vary from one set of users to another as well as language to language, and which uses are best facilitated by which combination of transcription efforts --- both human and machine.} 

\end{document}
